% BHU PYQ -- Auto-generated & error-corrected
\documentclass[11pt,a4paper]{article}

\usepackage[a4paper,top=2.5cm,bottom=2.5cm,left=2.8cm,right=2.8cm,headheight=14pt]{geometry}
\usepackage{amsmath,amssymb}
\usepackage[T1]{fontenc}
\usepackage[utf8]{inputenc}
\usepackage{lmodern}
\usepackage{microtype}
\usepackage{fancyhdr}
\usepackage{enumitem}
\usepackage{hyperref}
\usepackage{lastpage}
\usepackage{parskip}

\hypersetup{colorlinks=true,urlcolor=black,linkcolor=black,
  pdftitle={STB-401: Sample Surveys and Design of Experiments -- BHU PYQ},pdfauthor={Banaras Hindu University}}

\pagestyle{fancy}\fancyhf{}
\renewcommand{\headrulewidth}{0.4pt}\renewcommand{\footrulewidth}{0.4pt}
\fancyhead[L]{\small   \textit{Banaras Hindu University}}
\fancyhead[R]{\small   \textit{STB-401: Sample Surveys and Design of Experi}}
\fancyfoot[C]{\small Page  \thepage\ of \pageref{LastPage}}


\newlist{parts}{enumerate}{2}
\setlist[parts,1]{label=(\alph*),leftmargin=*,itemsep=5pt,topsep=3pt}
\setlist[parts,2]{label=(\roman*),leftmargin=*,itemsep=3pt,topsep=2pt}

\newcommand{\pts}[1]{\hfill   \textit{[\#1]}}


\begin{document}


\begin{center}
  {\large\bfseries BANARAS HINDU UNIVERSITY}\\[2pt]
  {\normalsize B.Sc. (Hons.) Semester V Examination, 2023-24}\\[6pt]
  \rule{0.8\linewidth}{0.5pt}\\[6pt]
  {\large\bfseries Paper No. STB-401: Sample Surveys and Design of Experiments}\\[4pt]
  {\normalsize Time: 3 Hours\hfill Maximum Marks: 70}
\end{center}
\rule{\linewidth}{0.4pt}
\smallskip



\noindent   \textit{Instructions: Answer   \textbf{five} questions including Question~1 which is compulsory. Symbols have their usual meanings.}
\bigskip
\medskip


\noindent   \textbf{Question 1.} Write a brief note on the estimation of sample size in the case of simple random sampling, addressing any two methods known to you. For a population of size $N = 430$, we roughly know that $ar{Y} = 19$ and $S^2 = 85.6$. With SRS, what should be the size of the sample to estimate $ar{Y}$ with a margin of error of $10\%$ of $ar{Y}$, apart from a $1$ in $20$ chance? \pts{14}

\medskip\rule{\linewidth}{0.2pt}

\medskip


\noindent   \textbf{Question 2.} What do you understand by error in sampling? Discuss, in detail, the different types of errors involved in sampling. \pts{14}

\medskip\rule{\linewidth}{0.2pt}

\medskip


\noindent   \textbf{Question 3.} Explain the method of optimum allocation of sample sizes for different strata when the total sample size $n$ is fixed. A sampler proposes to take a stratified random sample. The field costs are expected to be of the form $\sum c_i n_i$ and the advance estimates of relevant quantities for the two strata are as follows:

\begin{center}

\begin{tabular}{|c|c|c|c|}
\hline
Stratum & $p_i = N_i/N$ & $S_i$ & $c_i$ \\
\hline
1 & 0.4 & 10 & 4 \\
2 & 0.6 & 20 & 9 \\
\hline
\end{tabular}
\end{center}
Considering that the notations have their usual meanings:

\begin{parts}
  \item Find the values of $n_1/n$ and $n_2/n$ that minimize the total cost for a given value of $V(ar{y}_{st})$.
  \item Find the sample size required, under this optimum allocation, to make $V(ar{y}_{st}) = 1$, if fpc is ignored.
  \item Obtain the total field cost. \pts{14}
\end{parts}

\medskip\rule{\linewidth}{0.2pt}

\medskip


\noindent   \textbf{Question 4.} Discuss the concept of systematic sampling with its merits, demerits, and uses. Obtain the expression for the variance of the estimated mean under systematic sampling. \pts{14}

\medskip\rule{\linewidth}{0.2pt}

\medskip


\noindent   \textbf{Question 5.} What do you mean by the regression method of estimation? Find the bias and variance of the regression estimator for the population mean. Also prove that the regression estimator for the population mean is more efficient than the ratio estimator. \pts{14}

\medskip\rule{\linewidth}{0.2pt}

\medskip


\noindent   \textbf{Question 6.} What do you understand by analysis of variance (ANOVA)? Mention the assumptions underlying the ANOVA test. Describe ANOVA for a fixed-effects model for two-way classified data with one observation per cell. \pts{14}

\medskip\rule{\linewidth}{0.2pt}

\medskip


\noindent   \textbf{Question 7.} Explain the basic principles of design of experiments with examples. Also discuss the layout of a Randomized Block Design along with its disadvantages. \pts{14}

\medskip\rule{\linewidth}{0.2pt}

\medskip


\noindent   \textbf{Question 8.} Write a note on the design which follows all the basic principles of design of experiments. Mention its advantages and disadvantages over other designs. \pts{14}

\vfill
\rule{\linewidth}{0.4pt}

\begin{center}\small
  \href{https://bhu-exam.vercel.app/}{Click here to visit our website}\\[4pt]\href{https://me-aryan.vercel.app/}{Click here to visit our contributor}\\[4pt]\href{https://quashan-me.vercel.app/}{Click here to visit our contributor}
\end{center}

\end{document}